\section{Benchmark data-generating process}\label{app:dgp}
Random fields are unit-variance, kernel-smoothed white noise, $g_\ell=\mathrm{std}(K_\ell\zeta)$ with $\zeta\sim\mathcal N(0,I)$ and $K_\ell(i,j)=\exp\{-d_{ij}^2/(2\ell^2)\}$, where $d_{ij}$ is the distance between district centroids in km; $\nu$ denotes independent standard normal draws.

\paragraph{Rainfall.} $\mathrm{rain}_{itm}=N_i\varphi_m\exp(z_{itm}-0.065)$ with $z_{itm}=-0.2\,e_t c_m+0.18A_{it}+0.30M_{itm}+0.08\nu_{itm}$, where $N_i$ is the rainfall normal, $\varphi_m=(0.11,\allowbreak 0.13,\allowbreak 0.14,\allowbreak 0.18,\allowbreak 0.20,\allowbreak 0.09,\allowbreak 0.02)$ for May--November, $c_m=1.4$ in June--July and 1 otherwise, and $A_{it}$ and $M_{itm}$ are fields with $\ell=150$ and 80~km.

\paragraph{Temperature and humidity.} With $\tilde z=\mathrm{clip}(z/0.36,-3,\allowbreak 3)$ and $\theta_{it}=0.30e_t+0.25B_{it}$ ($B$ a 150-km field): $T_{itm}=\bar T_m-0.0065(h_i-170)+0.035(t-2008)+\theta_{it}-0.35\tilde z_{itm}+0.15\nu$; $T^{\max}_{itm}=T_{itm}+6+0.8\max(0,-\tilde z_{itm})+0.3\nu$; and $\mathrm{RH}_{itm}=\mathrm{clip}(\overline{\mathrm{RH}}_m+4\tilde z_{itm}-1.5\theta_{it}+1.5\nu,\,40,\,98)$, where $h_i$ is elevation and $\bar T_m$, $\overline{\mathrm{RH}}_m$ are regional monthly climatologies.

\paragraph{Soil moisture and reservoirs.} $s_{itm}=0.45\,s_{it,m-1}+0.55\{0.08+0.42(1-e^{-\mathrm{rain}_{itm}/200})(1.30-\mathrm{sand}_i/100)\}$; the observed series adds $\mathcal N(0,\allowbreak 0.025^2)$ error. Storage at sowing is $\mathrm{clip}(58-14e_{t-1}+\delta_b+7\nu,\,12,\,95)$\% with $\delta_b=-8$, 0 and 6 for the Mekong, Mun and Chi basins, and the irrigation supply is $S_{it}=I_i\times\mathrm{storage}_{b(i)t}/100$.

\paragraph{Prices.} $\ln P_t=\mu+0.6(\ln P_{t-1}-\mu)+\Sigma^{1/2}\nu_t$ plus occasional known shocks, with $\mu=\ln(10.5,\allowbreak 2.2,\allowbreak 1.0,\allowbreak 8.0,\allowbreak 14.0)$~THB\,kg$^{-1}$ for rice, cassava, sugarcane, maize and urea, standard deviations $(0.11,\allowbreak 0.17,\allowbreak 0.11,\allowbreak 0.13,\allowbreak 0.13)$, and correlations of 0.35 among crops and 0.2 with urea. Local prices are $P_{c,t-1}-\kappa_c d_{ic}$ with $\kappa=(0.010,\allowbreak 0.006,\allowbreak 0.004,\allowbreak 0.010)$~THB\,kg$^{-1}$\,km$^{-1}$ and the distance to a market or, for cane, to a sugar mill. The competing-revenue index is
\[ c_{it}=\ln\frac{p_{\mathrm{cas}}\bar Y_{\mathrm{cas}}+w_i\,p_{\mathrm{cane}}\bar Y_{\mathrm{cane}}+p_{\mathrm{mz}}\bar Y_{\mathrm{mz}}}{(2+w_i)\,p_{\mathrm{rice}}\bar Y_{\mathrm{rice}}},\] with $w_i=e^{-d^{\mathrm{mill}}_i/40}$ and typical yields $\bar Y=(400,\allowbreak 3500,\allowbreak 11000,\allowbreak 700)$~\kgrai{} for rice, cassava, sugarcane and maize; the input index is $\ln(p_{\mathrm{rice}}/P_{\mathrm{urea}})$. Both indices are standardised before entering Table~\ref{tab:dgp}. After each season the rice price is multiplied by $\exp(-\kappa\bar a_t/100)$, where $\bar a_t$ is the regional mean yield deviation from the base level and trend.

\paragraph{Satellite indices.} With visible yield $v_{it}$ (all components except noise), $\mathrm{NDVI}^{\mathrm{rice}}_{itm}=0.20+(N^{\mathrm{pk}}_{it}-0.20)\exp\{-((m-m^{\mathrm{pk}}_{it})/1.7)^2\}$, $N^{\mathrm{pk}}_{it}=0.42+0.38/\{1+\exp(-(v_{it}-380)/85)\}$ and $m^{\mathrm{pk}}_{it}=9.4+0.8\,\mathrm{clip}((0.25-\bar s^{MJ}_{it})/0.05,\allowbreak 0,\allowbreak 1.5)$; flooding reduces September--October values by up to 25\%, and observed NDVI is $0.85\,\mathrm{NDVI}^{\mathrm{rice}}+0.15\times0.62+\mathcal N(0,\allowbreak 0.025^2)$. EVI equals $0.12+(\mathrm{EVI}^{\mathrm{pk}}-0.12)$ times the same seasonal shape plus $\mathcal N(0,\allowbreak 0.02^2)$, with $\mathrm{EVI}^{\mathrm{pk}}=0.22+0.00075(v-250)$; LSWI is $-0.10+0.9(s-0.2)$ plus transplanting and flood terms and $\mathcal N(0,\allowbreak 0.03^2)$; LST is $T^{\max}+1.5+25(0.28-s)-6(\mathrm{NDVI}-0.5)+\mathcal N(0,\allowbreak 0.8^2)$. With monthly cloud probabilities $p_m=(0.15,\allowbreak 0.35,\allowbreak 0.45,\allowbreak 0.50,\allowbreak 0.45,\allowbreak 0.25,\allowbreak 0.10)$, a composite is lost with probability $0.25p_m$ and contaminated with probability $0.75p_m$ (NDVI reduced by $U(0.05,\allowbreak 0.18)$, EVI by $U(0.03,\allowbreak 0.10)$, LST by $U(2,\allowbreak 6)$\,$^\circ$C, LSWI increased by $U(0,\allowbreak 0.08)$).

\section{Additional results}\label{app:extra}
This appendix documents the tuning of LightGBM and gives the detailed numbers behind Fig.~\ref{fig:seasons}, Fig.~\ref{fig:attribution} and Table~\ref{tab:robust}. Table~\ref{tab:hparams} lists the random-search grid, the configuration selected on the validation seasons and the validation RMSE of the three candidate main models. Table~\ref{tab:byseason} gives the test RMSE of every learner in every season: the failures of the linear learners (Ridge in \worstYearRidge{} and the panel fixed-effects model in \worstYearFE) stand out against otherwise competitive errors, and outside these seasons the ensemble is usually the most accurate learner. Table~\ref{tab:ablation} reports the drop-one-group ablation for all three learners: for LightGBM only the satellite group matters, whereas for Ridge and the ensemble the removal of most groups \emph{lowers} the error---a symptom of the extrapolation failures diagnosed in Section~\ref{sec:diag}. Table~\ref{tab:worlds} lists every replicate and scenario world; the failures of Ridge, and of the ensemble that contains it, are confined to a few worlds, whereas LightGBM's error varies little.
\begin{table}[H]\centering\footnotesize
\caption{LightGBM random search (15 configurations $\times$ 2 target formulations, scored by rolling origin on 2015--2017) and the selected configuration (target: level). Validation RMSE of the candidate main models (kg\,rai$^{-1}$): Ensemble (LightGBM + Ridge) 26.2 ; Ridge regression 27.4 ; LightGBM 28.1.}\label{tab:hparams}
\begin{tabular}{@{}llc@{}}
\toprule
Hyper-parameter & Search grid & Selected \\
\midrule
\texttt{num\_leaves} & 7, 15, 31, 63 & 63 \\
\texttt{learning\_rate} & 0.01, 0.02, 0.03, 0.05, 0.08 & 0.05 \\
\texttt{n\_estimators} & 200, 400, 700, 1000 & 400 \\
\texttt{min\_child\_samples} & 10, 20, 40, 80 & 80 \\
\texttt{subsample} & 0.6, 0.8, 1.0 & 0.8 \\
\texttt{colsample\_bytree} & 0.3, 0.5, 0.7, 0.9 & 0.5 \\
\texttt{reg\_lambda} & 0.0, 1.0, 5.0, 20.0 & 1 \\
\texttt{max\_depth} & -1, 4, 6, 8 & -1 \\
\bottomrule
\end{tabular}
\end{table}
\begin{table}[H]\centering\footnotesize
\caption{Test RMSE (kg\,rai$^{-1}$) of every learner in every test season (bold: lowest error of the season) and over all 1{,}600 forecasts.}\label{tab:byseason}
\begin{tabular}{@{}lccccccccc@{}}
\toprule
Model & 2018 & 2019 & 2020 & 2021 & 2022 & 2023 & 2024 & 2025 & All \\
\midrule
LightGBM & 28.7 & 27.0 & 27.0 & 26.4 & 26.1 & 27.1 & 31.9 & 24.4 & 27.4 \\
Random forest & 27.4 & 25.7 & 28.9 & 28.6 & 26.9 & 29.2 & 32.3 & 25.9 & 28.2 \\
Ensemble (LightGBM + Ridge) & 50.9 & \textbf{24.8} & \textbf{26.3} & \textbf{25.6} & \textbf{24.2} & 24.3 & 29.5 & \textbf{23.7} & 29.9 \\
LightGBM (alternative target) & 28.0 & 34.0 & 28.5 & 29.0 & 28.0 & 33.7 & 36.7 & 30.2 & 31.2 \\
Multilayer perceptron & 35.8 & 30.3 & 35.9 & 29.7 & 33.7 & 30.3 & 33.0 & 27.8 & 32.2 \\
Panel fixed effects & \textbf{27.1} & 27.9 & 65.4 & 27.1 & 27.3 & 29.0 & 34.7 & 24.0 & 35.2 \\
Ridge regression & 81.7 & 25.3 & 27.2 & 28.2 & 25.2 & \textbf{24.3} & \textbf{28.5} & 24.8 & 37.9 \\
Climatology (3-season mean) & 29.9 & 49.3 & 43.4 & 54.5 & 42.4 & 62.6 & 47.2 & 39.8 & 47.1 \\
Persistence & 44.8 & 51.1 & 64.0 & 58.0 & 37.2 & 71.1 & 84.1 & 50.1 & 59.2 \\
\bottomrule
\end{tabular}
\end{table}
\begin{table}[H]\centering\scriptsize
\caption{Drop-one-group ablation: change in test RMSE (\%) when a feature group is removed and the learner re-trained in every fold; brackets: 95\% season-block bootstrap interval. Positive values mean the group is useful. The first row gives the RMSE (kg\,rai$^{-1}$) with all features.}\label{tab:ablation}
\begin{tabular}{@{}lccc@{}}
\toprule
Removed group & LightGBM & Ridge & Ensemble \\
\midrule
(none: RMSE with all features) & 27.4 & 37.9 & 29.9 \\
\midrule
Satellite (NDVI, EVI, LSWI, LST) & +25.0 [+13.0, +41.3] & +153.3 [+12.9, +279.8] & +90.3 [+14.1, +170.6] \\
Map layers \& spatial lags & +1.6 [-0.5, +3.7] & -22.2 [-67.5, +12.9] & -9.5 [-34.0, +6.9] \\
Geography (terrain, soil) & +1.1 [-0.7, +2.7] & +30.0 [-0.4, +60.0] & +15.1 [-0.5, +34.3] \\
Prices (own \& competing crops) & +0.4 [-1.3, +2.2] & -31.4 [-74.4, +0.1] & -13.9 [-37.9, +1.3] \\
Humidity \& soil moisture & -0.6 [-2.9, +1.8] & -30.4 [-74.4, +1.9] & -14.1 [-39.3, +1.6] \\
Climate (temperature, GDD) & -0.7 [-2.3, +1.0] & -28.5 [-70.1, +2.3] & -13.2 [-36.0, +1.4] \\
Water (rain, reservoir, irrigation) & -0.9 [-3.1, +1.0] & -8.3 [-20.3, +3.8] & -4.7 [-13.1, +2.3] \\
Previous yields & -2.8 [-4.8, -0.7] & -16.3 [-42.4, +6.5] & -8.2 [-22.8, +1.9] \\
\bottomrule
\end{tabular}
\end{table}
\begin{table}[H]\centering\footnotesize
\caption{Test RMSE (kg\,rai$^{-1}$) in every replicate world (seeds 1000--1009) and scenario world (seed 42), using the configuration tuned once in the main world. $^{\dagger}$: RMSE more than 1.5 times that of LightGBM in the same world. No sat.: ensemble without satellite features; Random CV: the same ensemble scored by random 5-fold cross-validation.}\label{tab:worlds}
\begin{tabular}{@{}lcccccc@{}}
\toprule
World & Climatology & Ridge & LightGBM & Ensemble & No sat. & Random CV \\
\midrule
Replicate world 1000 & 46.4 & 26.2 & 26.2 & 24.9 & 31.4 & 23.4 \\
Replicate world 1001 & 48.1 & 26.0 & 25.6 & 24.4 & 31.0 & 23.4 \\
Replicate world 1002 & 52.4 & 91.7$^{\dagger}$ & 26.4 & 50.3$^{\dagger}$ & 70.2 & 23.3 \\
Replicate world 1003 & 46.8 & 27.3 & 27.1 & 25.9 & 30.1 & 24.4 \\
Replicate world 1004 & 49.2 & 29.7 & 26.3 & 25.6 & 38.0 & 23.6 \\
Replicate world 1005 & 54.2 & 34.3 & 27.1 & 28.4 & 39.2 & 23.5 \\
Replicate world 1006 & 47.9 & 51.8$^{\dagger}$ & 27.0 & 33.9 & 48.4 & 23.7 \\
Replicate world 1007 & 53.4 & 28.1 & 26.8 & 24.9 & 29.3 & 23.1 \\
Replicate world 1008 & 48.1 & 100.7$^{\dagger}$ & 26.2 & 55.2$^{\dagger}$ & 35.6 & 23.1 \\
Replicate world 1009 & 47.0 & 26.1 & 26.1 & 25.1 & 31.6 & 23.6 \\
\midrule
Stronger non-linearity & 87.2 & 157.5$^{\dagger}$ & 28.5 & 83.9$^{\dagger}$ & 166.8 & 26.2 \\
3$\times$ map-layer error & 47.1 & 39.3 & 27.6 & 30.5 & 55.4 & 23.8 \\
1.7$\times$ cloud contamination & 47.1 & 54.7$^{\dagger}$ & 28.6 & 36.8 & 65.1 & 24.6 \\
No satellite-only pests & 46.8 & 25.8 & 26.3 & 25.0 & 43.6 & 23.6 \\
\bottomrule
\end{tabular}
\end{table}
